\section{Clinical Properties and Fuzzy Allocation}
\label{sec:properties_allocation}

\subsection{Stage three, the six clinical properties}
\label{sec:predicates}
Each of the six properties
names a way a denoised OCT image is known to fail and each is written as a formula so it can be
recomputed by anyone. Table~\ref{tab:predicates} gives, for every property, the reason it exists,
its score $s_i \in [0,1]$ as a weighted sum of named subterms, its map
$\mathbf{m}_i \in [0,1]^{H \times W}$ in closed form, where that map is large, and the threshold
$\tau_i$ that separates satisfied from unsatisfied, with every subterm defined beneath the table. A
larger score means the property is better satisfied and a larger map value means stronger local
evidence that it is not. None of these quantities needs a clean reference. They are computed from the
noisy input and the backbone output only, so they remain available at test time on a scanner the
method has never seen.

\subsection{Stage four, the rule layer}
\label{sec:negotiator}
The rule layer decides where the gain proposal should act through five named rules, each casting a
signed vote.

Each property score is first turned into a soft failure signal by comparing it with its threshold,
\begin{equation}
f_i = \sigma\big(\kappa (\tau_i - s_i)\big), \qquad \kappa = 10 ,
\label{eq:soft_failure}
\end{equation}
so $f_i$ is near one when property $i$ is failing and near zero when it is satisfied. The same soft comparison is used to test the cooperation
map and the strength of a proposal against their own thresholds. Each corrector produces that
strength as a per-pixel map from the backbone output and the failure maps it serves, large where its
property is far below threshold. Before the rules fire, that strength is multiplied by one half plus the
demand map, so the same proposal looks stronger to the rules at a pixel where the head reports low
cooperation. Conjunction uses the Lukasiewicz
form $\mathrm{AND}(a,b) = \max(0, a + b - 1)$, which is piecewise linear and has a bounded slope,
which is what makes the bound in Section~\ref{sec:stability} tight.

The five rules are the following. Rule one fires when the cooperation map is high and the proposal is
weak, and it votes against changing the image. Rule two fires when the map is low and the
proposal is strong, and it votes for changing it. Rule three fires when a relevant property is
failing, and it votes for changing it. Rule four fires when the map and the proposal strength are both
middling, and it votes weakly for changing it. Rule five fires when several correctors want the same
pixel, and it votes against changing it. Write $r_1$ to $r_5$ for the resulting truth values.

The allocation at each pixel is the bounded weighted vote
\begin{equation}
\begin{aligned}
a ={}& \underbrace{B_{\min} + (B_{\max} - B_{\min})\,\sigma(\beta)}_{\text{base, no rule firing}} \\
    &- C_1 w_1 r_1 + C_2 w_2 r_2 + C_3 w_3 r_3 + C_4 w_4 r_4 - C_5 w_5 r_5 ,
\end{aligned}
\label{eq:allocation}
\end{equation}
where $\beta$ is a learnable base logit, $w_k = \sigma(\theta_k) \in (0,1)$ are learnable rule weights
and the budgets are $B_{\min} = 0.25$, $B_{\max} = 0.50$, $C_1 = 0.15$, $C_2 = 0.25$, $C_3 = 0.15$,
$C_4 = 0.10$ and $C_5 = 0.10$.

The budgets satisfy $B_{\max} + C_2 + C_3 + C_4 = 1$ and
$B_{\min} - C_1 - C_5 = 0$, so since every truth value lies in $[0,1]$ and every weight in $(0,1)$,
the value of Equation~\eqref{eq:allocation} lies in $[0,1]$ for every input and every parameter value
and no clipping is ever needed. Without the identity a clamp is needed, the allocation saturates at
one end of it, as it did on every pixel in the submitted version of this work, the rule weights
receive no gradient, and the stability result of Section~\ref{sec:stability} is vacuous.
\label{sec:allocation_defect}
